\section{Related work}
\label{sec:related_target_centering}

Physical knowledge enters learned models through several interfaces. Hybrid
process models combine mechanistic structure with learned quantities~\cite{PsichogiosUngar1992};
physics-informed neural networks use differential-equation residuals in the
objective~\cite{RaissiPerdikarisKarniadakis2019}; APHYNITY learns a constrained
correction to incomplete dynamics~\cite{YinEtAl2021}. These approaches motivate
studying which part of a prediction problem is assigned to the physical
description and which part remains statistical. Here the physical prescription
defines the response coordinate, and the flow learns its conditional density.

Residual targets have a direct precedent in $\Delta$-machine learning, where
a learned correction augments an inexpensive quantum-chemical baseline~\cite{RamakrishnanEtAl2015}.
Conditional transformation models extend response transformation to
distributional regression by estimating transformations that depend on
explanatory variables~\cite{HothornKneibBuhlmann2014}. These lines connect
baseline correction with conditional distribution modeling. Our center is
prescribed before flow fitting, so the experimental variable is the residual
task presented to a matched density estimator.

Response parameterization can affect both model adequacy and optimization.
Box and Cox studied transformations that make a simple regression and error
model appropriate, comparing likelihoods on the original response scale with
the Jacobian included~\cite{BoxCox1964}. PopArt changes target normalization
while compensating the output layer to preserve the unnormalized predictor,
relating normalization to modified squared-loss SGD updates~\cite{VanHasseltEtAl2016}.
For squared-loss kernel regression, target power across kernel eigenfunctions
links target--model alignment to generalization~\cite{CanatarBordelonPehlevan2021}.
Together these results show why an information-preserving coordinate choice
can still matter to a restricted learner. Our arms share a global target scale
and fit a residual offset, focusing the comparison on input-dependent centering.

Normalizing flows permit tractable density evaluation through change of
variables. Conditional masked autoregressive flows incorporate explanatory
inputs into their transformations~\cite{PapamakariosPavlakouMurray2017},
and neural spline flows provide monotone rational-quadratic maps with analytic
inverses~\cite{DurkanEtAl2019}. Tail fidelity and expressiveness also depend on
base-distribution structure, Lipschitz behavior, depth and numerical
conditioning~\cite{JainiKobyzevYuBrubaker2020,KoehlerMehtaRisteski2021}.
Numerical conditioning concerns the invertible map and its Jacobian, distinct
from the physical information supplied to the neural conditioner.

In high-energy physics, deep generative surrogates increasingly complement
detailed Monte Carlo detector transport~\cite{HashemiKrause2024}, with
conditional normalizing flows achieving high-fidelity shower and material-parameterized
response modeling~\cite{KrauseShih2023,ErdmannEtAl2025}. Conversely, in inverse
detector reconstruction, physics-constrained generative models have recently
recovered sub-crystal electromagnetic shower structure from coarse segmented
calorimeter readouts by enforcing low-order spatial-moment
constraints~\cite{LiuTung2026}. Charged-particle tracking frameworks likewise
combine fast analytic material-interaction formulas with stochastic straggling
descriptions~\cite{AiEtAl2022}. In such forward hybrid settings, an analytic
stopping-power reference can evaluate exact step kinematics directly, whereas a
learned residual model may operate on coarsened, tabulated or reduced
conditioning inputs.

These studies establish residual modeling, coordinate choice and flow-based
transport surrogates as relevant tools; they do not determine the ordering of
finite spline-flow NLLs under restricted neural inputs. Our contribution is
the controlled information-axis comparison for a prescribed physical target
center. Matched raw-target baselines separate centering benefit from the
absolute benefit of additional neural information, and regional evaluations
test how aggregate changes are distributed across the physical domain.
